\makeabstract
    {
    Synthesis of Naphthalenediimide Threonine Derivatives for Nanotube Formation
    }
    {
    Annika Xuan\textsuperscript{1}, Makafui Baragbor\textsuperscript{1}, David E. Hansen\textsuperscript{1}
    }
    {
    \textsuperscript{1} Department of Chemistry, Amherst College, Amherst, MA
    }
    {
    Supramolecular chemistry investigates noncovalent forces that promote the self-assembly of larger structures, often to model biological systems such as proteins. As observed by the Sanders group, Naphthalenediimide (NDI) derivatives form nanotube structures through hydrogen bonding. Nanotube stability can then be increased by selecting for amino acid hydrophobicity and by adding hydrophobic covalent linkers. 
    }
    {
    Naphthalenediimide derivatives were synthesized using microwave synthesis to unsymmetrically add L-Phenylalanine and L-Threonine to the naphthalene core. Coupling procedures were adapted using EDC coupling and acylation techniques to yield a dimer NDI with a terephthalic ester linker. These reactions were adjusted due to the EPA ban on dichloromethane, with dimethyl carbonate and acetonitrile tested as possible solvent replacements. The resulting compounds were observed and characterized using Nuclear Magnetic Resonance (NMR).
    }
    {
    EDC coupling using dimethyl carbonate yielded the greatest amounts of mono-coupled and di-coupled material, but a large amount of starting material remained. The dimer product was identified by the \ensuremath{\beta} proton shift upfield in the NMR spectrum. When run at higher temperatures, EDC coupling produced an elimination byproduct in greater yield than the dimer. Attempts at acylation of L-threonine alone were primarily successful in producing the mono-acylated material but were abandoned because TFA, which is used in a subsequent procedure to remove the t-butyl protecting groups, was required. 
    }
    {
    EDC coupling of the NDI L-threonine derivatives was identified as the most suitable procedure for dimerization. These findings will allow for comparison with the NDI L-serine derivative counterpart to examine the effects of the methyl group on nanotube stability and hydrophobicity. Future directions toward this analysis require a procedure to deprotect the t-butyl protecting groups and CD spectra characterization. Additional tuning towards the covalent linker hydrophobicity can also be conducted using naphthalene and anthracene linkers. 
    }

    \newpage
    